\begin{abstract}
The double descent phenomenon---wherein test error rises near the interpolation threshold before falling again---challenges classical bias-variance intuitions about over-parameterization. While data augmentation is widely known to improve generalization, its mechanistic effect on the double descent curve has received little systematic study. We train ResNet and MLP models across seven parameter scales (11K--2.78M) with five augmentation strategies (None, Flip+Crop, Cutout, MixUp, Color Jitter) on CIFAR-10, repeated over three random seeds. We find that \textbf{ResNets without augmentation exhibit a clear double descent}, with test error peaking at $\sim$697K parameters (0.1144) before rising further to 0.1331 at 2.78M. All four augmentation strategies completely eliminate this peak, yielding monotonically decreasing error curves with test error as low as 0.0407 at 2.78M (MixUp). In contrast, \textbf{MLPs show no double descent} regardless of augmentation, suggesting this phenomenon is specific to convolutional architectures. A MixUp intensity ablation reveals negligible sensitivity to $\alpha \in [0.1, 1.0]$ (differences $<$0.0001), a negative result that constrains mechanistic accounts. Experiments on CIFAR-100 confirm that augmentation benefits scale with task difficulty, with Cutout reducing error by 36.7\% relative to no augmentation at 2.78M parameters.
\end{abstract}
